\chapter{Археология исходной идеи: от отдельного объекта к категории редукций}

После неудачи абсолютной меры перечитываются ранние философские основания
проекта. В первом томе фундаментальной сущностью назывался не изолированный
объект, а закон допустимых переходов между описаниями --- категория
редукций. Позднейшая работа подробно развила операторные части, но часто
заменяла закон сборки поиском одного абсолютного скалярного действия.

\section{Устойчивый образец проекта}

Наиболее устойчивыми оказались относительные и композиционные конструкции:
отношения определителей, голономии, конусы отображения, относительные нормы,
детерминантные линии, индексы дефектов и остаточные циклы. Повторяющиеся
неудачи относятся к абсолютным весам топологий, абсолютному масштабу,
односледовым действиям и общей нормировке.

\section{Литературная граница}

Неограниченные KK-соответствия, относительные спектральные тройки, склейка
BV--BFV, детерминантные и пфаффиановые линии и относительные модулярные
операторы являются известными структурами. Новая задача проекта состоит не
в их переименовании, а в построении конкретного замкнутого треугольника
редукций RPFT/UGSM/TOE.

\section{Относительная архитектура редукций}

\begin{definition}[Категория редукций]
Объектами служат геометрические, спектральные, корреляционные и граничные
реализации. Морфизмами являются соответствия, относительные спектральные
тройки или переносы суперсвязностей. Композиция задаётся склейкой или
произведением типа Каспарова, а вес --- относительным коциклом либо
функтором со значениями в детерминантной линии.
\end{definition}

При локальном изменении нормировки \(B(X)\) вес ребра меняется как
\begin{equation}
 \Gamma(X\to Y)\mapsto\Gamma(X\to Y)+B(Y)-B(X).
\end{equation}
Открытый путь сохраняет неоднозначность концов, а сумма по замкнутому циклу
инвариантна. Машинная проверка подтверждает аддитивную и мультипликативную
формы этого утверждения.

\begin{theorem}[Неоднозначность открытого ребра и инвариантность цикла]
Локальная перенормировка изменяет отдельный относительный переход на
кограницу, но её вклад сокращается в замкнутой композиции. Следовательно,
только замкнутый цикл может нести независящий от локальной схемы остаток.
\end{theorem}

\section{Треугольник трёх редукций}

Минимальная схема проекта имеет вид
\begin{equation}
 R_g\xrightarrow{T_{gs}}R_s\xrightarrow{T_{sc}}R_c
 \xrightarrow{T_{cg}}R_g.
\end{equation}
Переход от геометрии к спектру задан частично; переход от спектра к
корреляции задаётся тепловым оператором и логарифмом; обратный переход от
корреляции к геометрии отсутствует. Поэтому коцикл полного треугольника пока
не определён.

\begin{nogocriterion}[Запрет категориального переименования]
Само введение слов «категория», «функтор» и «коцикл» не считается новой
архитектурой. Необходимо предъявить нетривиальный морфизм \(T_{cg}\), его
область определения, закон композиции и воспроизводимый остаток замкнутого
цикла.
\end{nogocriterion}

\section{Вывод}

Относительная гипотеза допускается как ограниченное направление
исследования: она объясняет общую картину неудач и выделяет точное
препятствие, но ещё не даёт физики. Следующая проверка ---
\path{version5_reduction_triangle_cocycle_gate} --- должна либо построить
нетривиальный обратный морфизм, либо закрыть минимальный треугольник.

Машинный сертификат сохранён в
\path{s2t_v5_foundational_relative_architecture_gate_results.json}.